\section{Methods \label{sec:methods}}

This section describes the machine learning methods used to infer physical launch and environmental parameters from simulated projectile trajectories.
Two complementary model classes are considered: a multilayer perceptron (MLP) applied to manually engineered scalar trajectory descriptors, and gated recurrent unit (GRU) networks applied directly to time-resolved trajectory data.
The overall motivation is to compare a compact feature-based representation against a sequence-based representation that can, in principle, learn relevant temporal structure directly from the observed motion.

\subsection{Feature-Based Regression with a Multilayer Perceptron}

As a non-recurrent baseline, we use a multilayer perceptron that maps a fixed-dimensional vector of scalar trajectory features to the desired physical target parameters.
In the implementation, the MLP first flattens the input, then applies a sequence of fully connected linear layers with ReLU activations, followed by a final linear output layer.
For an input feature vector $\vb{x}\in\mathbb{R}^{d_{\mathrm{in}}}$, the model can be written schematically as
\begin{align*}
    \vb{h}_1 &= \sigma\left(W_1 \vb{x} + \vb{b}_1\right), \\
    \vb{h}_{\ell+1} &= \sigma\left(W_{\ell+1}\vb{h}_{\ell}+\vb{b}_{\ell+1}\right), \qquad \ell = 1,\dots,L-1, \\
    \hat{\vb{y}} &= W_{\mathrm{out}}\vb{h}_{L} + \vb{b}_{\mathrm{out}},
\end{align*}
where $\sigma$ denotes the ReLU\cite{nair2010rectified} activation function and $\hat{\vb{y}}$ is the predicted target vector.
The MLP experiments use scalar inputs with $24$ engineered features, and the output dimension is either $3$ for predicting the Cartesian initial velocity or $11$ for predicting the full set of selected launch and environmental parameters.

The scalar features are designed to summarize the most important geometric and kinematic properties of the trajectory.
They include the initial and final positions $(x_0,y_0,z_0)$ and $(x_T,y_T,z_T)$, the total flight time $T$, the displacement components $(\Delta x,\Delta y,\Delta z)$, the Euclidean start-to-end distance, the total arc length of the trajectory, the apex position $(x_{\mathrm{peak}},y_{\mathrm{peak}},z_{\mathrm{peak}})$, the time of the apex $t_{\mathrm{peak}}$, and finite-difference estimates of the initial and final velocity components and speeds.
These quantities provide physically interpretable information about range, curvature, maximum height, flight duration, and endpoint behavior.
For example, the initial finite-difference velocity estimate is computed from the first two sampled trajectory points, while the final velocity estimate is computed from the last two sampled points.
The total trajectory length is computed by summing the Euclidean distances between successive sampled positions.
The apex is identified as the point of maximal height $z$, which gives both the peak position and the corresponding time.

For some experiments, the desired target is not the spherical launch representation $(v_0,\theta,\phi)$, but the Cartesian initial velocity $(v_x,v_y,v_z)$.
This target is computed using
\begin{equation*}
    \pmqty{v_x\\v_y\\v_z} = v_0\pmqty{\sin\theta\cos\phi\\\sin\theta\sin\phi\\\cos\theta}
\end{equation*}
which is implemented as a preprocessing target transformation.
Using Cartesian velocity components avoids angular discontinuities in $\phi$ and gives a target representation that is often easier for regression models to learn.

\subsection{Sequence-Based Regression with Gated Recurrent Units}

The second model class is based on gated recurrent units, which are recurrent neural networks designed to process sequential data.
Instead of receiving a single engineered feature vector, the GRU receives a time series of trajectory observations such as
\begin{equation*}
    \vb{x}_1,\vb{x}_2,\dots,\vb{x}_T, \qquad \vb{x}_t = \begin{pmatrix}
        t_t & x_t & y_t & z_t
    \end{pmatrix}^{\top},
\end{equation*}
or an extended version that also contains velocity and acceleration estimates.
The GRU then returns a final hidden state, which is passed through a small feed-forward prediction head consisting of a linear layer, ReLU activation, dropout, and a final linear layer.
In the endpoint-augmented variants, a second prediction head is added to predict the two-dimensional hit point separately from the physical parameter vector.

A GRU updates its hidden state using gating mechanisms that control how much past information is retained and how much new information is incorporated.
For an input $\vb{x}_t$ and previous hidden state $\vb{h}_{t-1}$, the standard GRU equations are
\begin{align*}
    \vb{z}_t &= \sigma\left(W_z \vb{x}_t + U_z \vb{h}_{t-1} + \vb{b}_z\right), \\
    \vb{r}_t &= \sigma\left(W_r \vb{x}_t + U_r \vb{h}_{t-1} + \vb{b}_r\right), \\
    \tilde{\vb{h}}_t &= \tanh\left(W_h \vb{x}_t + U_h(\vb{r}_t \odot \vb{h}_{t-1}) + \vb{b}_h\right), \\
    \vb{h}_t &= (1-\vb{z}_t)\odot \vb{h}_{t-1} + \vb{z}_t \odot \tilde{\vb{h}}_t.
\end{align*}
Here, $\vb{z}_t$ is the update gate, $\vb{r}_t$ is the reset gate, $\tilde{\vb{h}}_t$ is the candidate hidden state, and $\odot$ denotes element-wise multiplication.\cite{cho2014gru}
The update gate allows the network to preserve information over many time steps, while the reset gate controls how strongly the previous hidden state contributes to the candidate state.
This makes GRUs suitable for time-series analysis because they learn a compressed representation of the trajectory history rather than treating each time point independently.
In this project, the final hidden state $\vb{h}_T$ is interpreted as a learned summary of the observed flight path and is used to predict the physical parameters.

The GRU models are particularly appropriate for this problem because projectile trajectories contain temporally ordered information.
Quantities such as curvature, deceleration, and asymmetry are not fully described by individual positions alone, but by their evolution over time.
A recurrent architecture can use early and late parts of the trajectory jointly, and it can also be evaluated on cropped sequences, where only a prefix of the trajectory is available.\cite{cho2014gru}
The training code supports random cropping of time-series batches by selecting only the first $k$ time points, where $k$ is sampled between fractions of the full sequence length.
This allows the model to learn from partial trajectories and is useful for studying how much of the flight path is needed to infer the underlying launch and environmental parameters.

\subsection{Data Preprocessing and Representation}

The raw simulator output consists of time arrays, three-dimensional positions, launch parameters, and metadata stored in immutable raw samples.
During preprocessing, each raw sample is converted into a processed representation containing trajectory channels $t,x,y,z$, scalar features, targets, and optionally metadata.
The dataset builder applies a configurable pipeline of transformations to each raw sample and then writes the selected features and targets into PyTorch tensors.
For time-series data, the resulting feature tensor has shape
\begin{equation*}
    (N, n_{\mathrm{features}}, n_{\mathrm{timepoints}}),
\end{equation*}
whereas scalar-feature data are stored as
\begin{equation*}
    (N, n_{\mathrm{features}}).
\end{equation*}
The target tensor has shape $(N,n_{\mathrm{targets}})$ in both cases.

A central preprocessing step for the GRU models is resampling.
Although the raw trajectories were generated with many time points, the GRU experiments use a reduced fixed-length representation with $100$ time points.
The resampling transform first chooses new time points between the minimum and maximum time of the trajectory and then interpolates the trajectory channels using a cubic spline.\cite{de1978practical}
The implemented sampling distributions include even, logarithmic, uniform, and Gaussian sampling over the normalized time interval.
In the main GRU training configurations, even resampling is used, so the trajectory is represented by uniformly spaced time samples.
This reduces computational cost while preserving a standardized input length for mini-batch training.

Additional time-series features are obtained by numerical differentiation.
Velocity channels are computed as
\begin{equation*}
    v_i(t) \approx \dv{r_i}{t}, \qquad i\in\{x,y,z\},
\end{equation*}
using numerical gradients of the position channels with respect to time.
Acceleration channels are then computed analogously from the velocity channels.
Some GRU configurations therefore use only $(t,x,y,z)$, while others use the extended feature set $(t,x,y,z,v_x,v_y,v_z,a_x,a_y,a_z)$.
The motivation for adding velocities and accelerations is that the forces in the physical simulation act through velocity-dependent drag and Magnus terms, so derivative information may make the underlying parameters easier to identify.
At the same time, numerical differentiation can amplify noise, so the benefit of these features must be evaluated empirically. \cite{chartrand2011numerical}

All models use standardization of features and targets. Standardization transforms each feature to have zero mean and unit variance, ensuring that all input features are on a comparable scale. This improves the stability and efficiency of gradient-based optimization, allowing neural networks to converge faster and preventing features with larger numerical values from disproportionately influencing the learning process.\cite{lecun2012efficient}
For scalar inputs, feature means and standard deviations are computed over the training samples, while for time-series inputs they are computed over both the sample and time dimensions for each feature channel.
Targets are standardized using the training-set mean and standard deviation for each output dimension.
The standardized quantities are therefore
\begin{equation*}
    \tilde{x} = \frac{x-\mu_x}{\sigma_x+\epsilon}, \qquad \tilde{y} = \frac{y-\mu_y}{\sigma_y+\epsilon},
\end{equation*}
with a small numerical constant $\epsilon=10^{-8}$ used for stability.
The same scaling factors are then applied to the corresponding test data, which avoids information leakage from the test set into the training process.
For later standalone evaluation runs, the saved feature and target scaling factors can be loaded and reused.
